\hypertarget{python-3.12-native-generics-pep-695-type-statement-and-subinterpreters}{%
\chapter{Python 3.12: Native Generics (PEP 695), Type statement, and
Subinterpreters}\label{python-3.12-native-generics-pep-695-type-statement-and-subinterpreters}}

\hypertarget{pep-695-type-parameter-syntax}{%
\subsection{18.1 PEP 695 Type Parameter
Syntax}\label{pep-695-type-parameter-syntax}}

Python 3.12 introduced a clean, native syntax for generic classes,
generic functions, and type aliases using type parameters enclosed in
square brackets.

\hypertarget{transition-from-legacy-declarations}{%
\subsubsection{1. Transition from Legacy
Declarations}\label{transition-from-legacy-declarations}}

Prior to Python 3.12, defining generic types required importing and
manually instantiating \passthrough{\lstinline!TypeVar!} variables,
along with explicitly defining variance constraints (covariant,
contravariant, or invariant):

\begin{lstlisting}[language=Python]
# Legacy Python 3.11-
from typing import TypeVar, Generic

T = TypeVar('T', covariant=True)  # Manual variance specification

class Stack(Generic[T]):
    pass
\end{lstlisting}

PEP 695 replaces this boiler-plate with a cleaner syntax:

\begin{lstlisting}[language=Python]
# Modern Python 3.12+
class Stack[T]:
    pass
\end{lstlisting}

Under this syntax: * CPython automatically creates the type variable
\passthrough{\lstinline!T!} (instance of
\passthrough{\lstinline!typing.TypeVar!}) and binds it to the class
scope. * \textbf{Static Auto-Variance}: Static type checkers (like
Mypy/Pyright) automatically infer the variance of
\passthrough{\lstinline!T!} by analyzing how the variable is used inside
the class definition, completely eliminating the need for manual
\passthrough{\lstinline!covariant=True!} or
\passthrough{\lstinline!contravariant=True!} flags.

\hypertarget{bound-and-constraint-specifications}{%
\subsubsection{2. Bound and Constraint
Specifications}\label{bound-and-constraint-specifications}}

Type parameter bounds and constraints are defined directly inside the
square brackets: * \textbf{Bound Syntax
(\passthrough{\lstinline!T: bound\_type!})}: Restricts the type variable
to a specific subclass or type.
\passthrough{\lstinline!python     def process[T: int](val: T) -> T:         return val!}
* \textbf{Constraint Syntax
(\passthrough{\lstinline!T: (type1, type2, ...)!})}: Restricts the type
variable to one of a set of explicit types.
\passthrough{\lstinline!python     def parse[T: (str, bytes)](data: T) -> T:         return data!}

\begin{center}\rule{0.5\linewidth}{0.5pt}\end{center}

\hypertarget{annotation-scopes-and-lazy-evaluation}{%
\subsection{18.2 Annotation Scopes and Lazy
Evaluation}\label{annotation-scopes-and-lazy-evaluation}}

In Python 3.11 and below, type parameter bounds and type aliases were
evaluated eagerly at module import time. This caused circular import
errors (if a type variable referenced a class defined later in another
module) and increased startup memory footprints. PEP 695 resolves this
by introducing \textbf{Annotation Scopes} and lazy evaluation.

\hypertarget{annotation-scopes}{%
\subsubsection{1. Annotation Scopes}\label{annotation-scopes}}

When the compiler encounters type parameters or the
\passthrough{\lstinline!type!} statement, it wraps their evaluation code
inside a new lexical scope called an \textbf{Annotation Scope}: * An
annotation scope is a nested compiler-generated scope (similar to a
hidden function block). * Variables and bounds defined inside this scope
are evaluated \textbf{lazily}only when they are explicitly queried at
runtime.

\hypertarget{bytecode-disassembly-of-the-type-statement}{%
\subsubsection{\texorpdfstring{2. Bytecode Disassembly of the
\texttt{type}
Statement}{2. Bytecode Disassembly of the type Statement}}\label{bytecode-disassembly-of-the-type-statement}}

Let's analyze how CPython compiles a modern type alias:

\begin{lstlisting}[language=Python]
type Vector3D[T] = tuple[T, T, T]
\end{lstlisting}

\hypertarget{compiled-bytecode}{%
\paragraph{Compiled Bytecode:}\label{compiled-bytecode}}

\begin{lstlisting}
  1           0 LOAD_CONST               0 ('Vector3D')
              2 LOAD_CONST               1 ('T')
              4 LOAD_CONST               2 (<code object Vector3D at 0x...>)
              6 MAKE_FUNCTION            0
              8 SET_FUNCTION_ATTRIBUTE   8 (closure/annotations helper)
             10 CALL_FUNCTION            0
             12 BUILD_TYPEALIAS
             14 STORE_NAME               0 (Vector3D)
\end{lstlisting}

\begin{itemize}
\tightlist
\item
  \textbf{\passthrough{\lstinline!MAKE\_FUNCTION!}}: Creates a lazy
  evaluation function from the nested code object representing the alias
  value \passthrough{\lstinline!tuple[T, T, T]!}.
\item
  \textbf{\passthrough{\lstinline!BUILD\_TYPEALIAS!}}: Pops the
  evaluation function, type parameters, and name, then constructs an
  instance of \passthrough{\lstinline!typing.TypeAliasType!}.
\item
  \textbf{\passthrough{\lstinline!STORE\_NAME!}}: Binds the type alias
  object to \passthrough{\lstinline!Vector3D!}.
\end{itemize}

\hypertarget{c-level-representation-of-typealiastype}{%
\subsubsection{\texorpdfstring{3. C-Level Representation of
\texttt{TypeAliasType}}{3. C-Level Representation of TypeAliasType}}\label{c-level-representation-of-typealiastype}}

At the C level, \passthrough{\lstinline!TypeAliasType!} is represented
by a dedicated struct wrapping the properties of the type alias: *
\passthrough{\lstinline!\_\_name\_\_!}: Name of the alias. *
\passthrough{\lstinline!\_\_type\_params\_\_!}: A tuple of type
parameters. * \passthrough{\lstinline!\_\_value\_\_!}: The aliased type.
It uses a C descriptor getter that evaluates the lazy function code
object on the first access, caching the result to avoid redundant
evaluations.

\begin{center}\rule{0.5\linewidth}{0.5pt}\end{center}

\hypertarget{pep-684-per-interpreter-gil-and-subinterpreters}{%
\subsection{18.3 PEP 684: Per-Interpreter GIL and
Subinterpreters}\label{pep-684-per-interpreter-gil-and-subinterpreters}}

CPython has supported subinterpreters via its C-API since Python 1.5.
However, because they all shared a single Global Interpreter Lock (GIL),
only one interpreter could execute bytecode at a time, preventing true
multi-core parallel execution. Python 3.12 introduced \textbf{PEP 684},
providing a dedicated, isolated GIL for each subinterpreter.

\hypertarget{c-level-structural-separation}{%
\subsubsection{1. C-Level Structural
Separation}\label{c-level-structural-separation}}

CPython separates runtime state into two major structures: *
\textbf{\passthrough{\lstinline!PyRuntimeState!}}
(\passthrough{\lstinline!pycore\_runtime.h!}): Process-wide, shared
global resources, including system memory allocators, signal handlers,
and GC arena lists. *
\textbf{\passthrough{\lstinline!PyInterpreterState!}}
(\passthrough{\lstinline!pycore\_pystate.h!}): Interpreter-specific
resources.

Under PEP 684, each \passthrough{\lstinline!PyInterpreterState!} is
allocated its own dedicated GIL struct
(\passthrough{\lstinline!struct \_gil\_runtime\_state!}). This allows
multiple interpreter instances to execute bytecode in parallel on
separate threads, leveraging separate CPU cores.

\begin{lstlisting}
CPython Process Layout with Per-Interpreter GIL:

======================= [ PyRuntimeState ] =======================
  /                                                             \
 [ PyInterpreterState A ]              [ PyInterpreterState B ]
    - sys.modules                         - sys.modules
    - Private Object Heap                 - Private Object Heap
    - Dedicated GIL A                     - Dedicated GIL B
        |                                     |
        v                                     v
   ( OS Thread 1 )                       ( OS Thread 2 )
==================================================================
\end{lstlisting}

\hypertarget{isolation-boundaries}{%
\subsubsection{2. Isolation Boundaries}\label{isolation-boundaries}}

To prevent concurrency race conditions, interpreters maintain strict
boundaries: * \textbf{Private Module Cache}: Each subinterpreter has its
own \passthrough{\lstinline!sys.modules!} dictionary. * \textbf{Heap
Isolation}: Python heap objects (\passthrough{\lstinline!PyObject!}
references) cannot be shared directly between interpreters. If
interpreter A attempts to read or write a
\passthrough{\lstinline!PyObject!} allocated in interpreter B: 1.
Reference counting checks
(\passthrough{\lstinline!Py\_INCREF!}/\passthrough{\lstinline!Py\_DECREF!})
will trigger race conditions. 2. Garbage collection passes running in
interpreter A will attempt to collect memory belonging to interpreter B,
causing memory corruption. * \textbf{Communication Channel}:
Subinterpreters communicate exclusively through serialized message
channels (such as \passthrough{\lstinline!\_xxsubinterpreters!}). Data
is serialized (e.g.~converted to raw bytes or using shared memory
structures via the C buffer protocol) and then re-materialized as fresh
objects in the target interpreter's private heap.

\hypertarget{parallel-execution-code-example}{%
\subsubsection{3. Parallel Execution Code
Example}\label{parallel-execution-code-example}}

We can spawn subinterpreters and execute code concurrently using the
\passthrough{\lstinline!\_xxsubinterpreters!} module:

\begin{lstlisting}[language=Python]
import _xxsubinterpreters as interpreters
import threading

def run_in_subinterpreter(interp_id, script):
    # Run script inside the isolated interpreter context
    interpreters.run_string(interp_id, script)

# Create a new subinterpreter with an isolated GIL
interp_id = interpreters.create()

script = """
import time
# Executes concurrently on a separate CPU core
result = sum(i * i for i in range(10_000_000))
print("Subinterpreter result calculation complete:", result)
"""

# Spawn a separate OS thread to run the subinterpreter in parallel
thread = threading.Thread(target=run_in_subinterpreter, args=(interp_id, script))
thread.start()
thread.join()

# Destroy the subinterpreter and release resources
interpreters.destroy(interp_id)
\end{lstlisting}

\begin{center}\rule{0.5\linewidth}{0.5pt}\end{center}
